\section{Two accounts of the Gospel's fruit}

The two interpretations give different answers to the question raised by the three measures. Chrysostom sees the disproportion between little leaven and much meal; his explanation encourages weak disciples who must live among those to whom they preach. Jerome gives accounts of the human person whose faculties are brought into agreement; Augustine's identification of leaven with love explains the affection that moves such a person to act. The first interpretation makes contact and example especially prominent. The second makes understanding and reformed desire especially prominent.

They also differ in how they develop the prayers. In the first interpretation, the requests for acceptable conduct, cleansing, protection and continued nourishment concern the life through which the community becomes an example. In the second, the same petitions describe the needs of a person whose knowledge, actions and desires must agree. The comparison does not divide public witness from private holiness: Paul describes transformed people whose conduct is known, and Jerome places the communication of Scripture within the Church.

Both require an account of hope that can survive present suffering. Chrysostom explains the Spirit's joy amid affliction; Augustine directs the sinner's waiting toward the resurrection already accomplished in Christ. Their different explanations leave the prayer from the depths intelligible after the promise of peace has been heard. The Mass can praise deliverance, ask for mercy and seek further nourishment within the same celebration. Its people have received gifts whose full effects they still await.
